\begin{frame}{혼동행렬: 네 개의 숫자}
  환자 100명 검사 $\to$ \kb{혼동행렬}(confusion matrix):
  TP=18, FP=2, FN=7, TN=73.
  \vskip0.4em
  \begin{center}
  \small
  \begin{tabular}{@{}lcc@{}}
    \toprule
    실제 \textbackslash\ 예측 & 양성 예측 & 음성 예측 \\
    \midrule
    실제 양성 & \textbf{True Positive (TP)} & \textbf{False
      Negative (FN)} \\
    실제 음성 & \textbf{False Positive (FP)} & \textbf{True
      Negative (TN)} \\
    \bottomrule
  \end{tabular}
  \end{center}
  \vskip0.6em
  \begin{itemize}
    \item 실제 양성 $= 18+7 = 25$ 명, 모델이 양성 예측 $= 18+2 = 20$ 명.
    \item \kb{Precision} $= \frac{TP}{TP+FP}$: ``양성이라고 예측한
      것 중 진짜 양성 비율'' --- \textbf{거짓 경보}를 얼마나 피했는가.
    \item \kb{Recall} $= \frac{TP}{TP+FN}$: ``실제 양성 중 잡아낸
      비율'' --- \textbf{놓친 양성}이 얼마나 적은가.
    \item \kb{F1} $= \frac{2 \cdot P \cdot R}{P + R}$:
      P와 R의 \textbf{조화평균} --- 둘 다 낮으면 F1도 낮아지도록,
      한쪽만 높다고 후하게 주지 않음.
  \end{itemize}
\end{frame}
